\section{Arithmetic bundles for the specified native hulls}
\label{sec:uniform-arithmetic-bundles}

The local equalities already proved in
Section~\ref{sec:general-tame-square-labels} admit an explicit realization
by arithmetic vector bundles on the actual level field. We construct
these bundles, compute their normalized degrees, and descend a
sufficiently powered weighted determinant to the field of moduli
$\Q$. This realizes ordinary arithmetic objects with the required
numerical normalization. It does not establish their membership in
the full possible-output family of IUT III, Corollary~3.12.

The relevant source conventions are the relative weighted determinant
and sufficiently divisible positive tensor powers in
\cite[Remark~3.9.5(vii)(Ob3)]{IUT3}, the label average in
\cite[Proposition~3.9(i)--(iii)]{IUT3}, and the normalized local
weights in \cite[Proposition~1.4(i), Remark~1.7.1]{IUT4}.
The additional Ind3 comparisons and the particular global
prime-strip identification in
\cite[Theorem~3.11(ii), Corollary~3.12]{IUT3} are separate requirements;
none is assumed in the arithmetic arguments below.

\subsection{The fixed native input and its reference}

Fix a member of the family in
Theorem~\ref{thm:powerfree-global-family}, with primes
$\ell\equiv43\pmod{60}$ and $p\equiv-1\pmod{30\ell}$, and write
\[
 D_A:\ y^2=x(x-A^2)(x+A^2+1),\qquad
 K=\Q(i,D_A[30\ell]).
\]
The field $K$ is Galois over $\Q$. At the specified place over $p$,
its completion is the field of \eqref{eq:gs-field-data}:
\[
 E=\Q_p(\mu_{30\ell},\varpi),\qquad
 \varpi^{15\ell}=b_0,\qquad
 b_0\in\Q_p,\qquad v_p(b_0)=2.
\]
Put
\begin{equation}\label{eq:ub-parameters}
 \begin{gathered}
 e=15\ell,\quad f=2,\quad d=ef,\quad
 h=(\ell-1)/2,\quad \kappa=(e-1)/e,\\
 \beta=\varpi^{(e+1)/2}/p,\qquad r_0=\varpi^{15},\\
 m_j=j+1,\quad N_j=d^{m_j-1},\quad D_j=d^{m_j},\\
 a_j=r_0^{j^2},\qquad
 k_j=\left\lfloor\frac{2j^2}{\ell}\right\rfloor+1,\\
 \eta_j=e k_j-(e-1)m_j\quad(1\le j\le h).
 \end{gathered}
\end{equation}
The element $\beta$ is a uniformizer by Lemma~\ref{lem:gs-field};
$\varpi$ has valuation $2/e$ and is not one. All roots are fixed
before the permitted arrows are chosen.

For each label define
\[
 \begin{gathered}
 T_j=E^{\otimes_{\Q_p}m_j},\qquad
 A_j=\OO_E^{\otimes_{\Z_p}m_j},\\
 B_j=\prod_{\operatorname{Gal}(E/\Q_p)^{m_j-1}}\OO_E
       \subset T_j.
 \end{gathered}
\]
The product decomposition uses the first factor as coefficient
field and has $N_j$ copies of $E$. In general $A_j\ne B_j$.
The native lattice is
$I=p^{-1}\log\OO_E^\times=\beta^{1-e}\OO_E$.
Let $\Gamma_E$ be the already constructed family of actual native
integral Kummer arrows on $I$, using the same arrow in repeated slots.
We retain precisely the marked point inputs
\begin{equation}\label{eq:ub-native-points}
 \begin{gathered}
 \rho=p^{-1}\log_{\rm BK}^{\rm std},\\
 u=\rho({\rm Kum}(1+p))=p^{-1}\log(1+p),\\
 \tau_j=\rho({\rm Kum}(1+pa_j))=p^{-1}\log(1+pa_j).
 \end{gathered}
\end{equation}
Their closed $B_j$-span after the permitted arrows is
\[
 P_j^\rho=
 \overline{\Span_{B_j}
   \{F(\tau_j)\otimes(Fu)^{\otimes j}:F\in\Gamma_E\}}.
\]
Corollaries~\ref{cor:gs-hulls} and~\ref{cor:gs-pilot-scale}, with the
per-arrow bridge of Proposition~\ref{prop:np-point-bridge}, give
\begin{equation}\label{eq:ub-local-hull}
 M_j^\rho:=\mathcal M_{\Gamma_E}(a_j)
       =P_j^\rho=\beta^{\eta_j}B_j.
\end{equation}
The first term is the pre-transport principal-ideal hull already
defined in \eqref{eq:np-preideal}. Equality here identifies these
specified native hulls, not their input points or the full published
theta-pilot family. No additional choice of a torsion representative,
vertical level, or Ind3 enlargement is included in
\eqref{eq:ub-native-points}--\eqref{eq:ub-local-hull}.

Let $\mu_{B_j}$ be ordinary additive Haar measure on $T_j$ with
$\mu_{B_j}(B_j)=1$, and write
\begin{equation}\label{eq:ub-local-volume}
 V_j^B(H)=D_j^{-1}\log\mu_{B_j}(H).
\end{equation}
The positive root $\mu_{B_j}^{1/D_j}$ is not another additive Haar
measure. For a general field of moduli $F_0$, the normalized weight
of an ordered word $(v_0,\ldots,v_j)$ over a rational prime is
\[
 \frac{\prod_{i=0}^j[(F_0)_{v_i}:\Q_p]}
      {[F_0:\Q]^{m_j}}.
\]
Indeed the sum of the numerator over all words is the $m_j$-th power
of $\sum_{v\mid p}[(F_0)_v:\Q_p]=[F_0:\Q]$. In the rational branch
there is one selected-place word and this weight is 1. The remaining
procession normalization is $h^{-1}\sum_{j=1}^h$, with no further
division by $m_j$.

\subsection{Actual fractional ideals and normalized degree}

At a finite place $w$ of $K$, put $q_w=\#k(w)$ and
$|x|_w=q_w^{-\operatorname{ord}_w(x)}$. This is the local norm
absolute value. At an infinite place $\sigma$, use ordinary real
or complex modulus and multiplicity $n_\sigma=1$ or $2$,
respectively. For a metrized line $\overline L$, declare its finite
local lattices to be the unit balls. For a nonzero rational section
$s\in L\otimes_{\OO_K}K$, define
\begin{equation}\label{eq:ub-degree}
 \widehat{\deg}_K(\overline L)
 =-\frac1{[K:\Q]}
   \left(\sum_{w\ {\rm finite}}\log\|s\|_w+
         \sum_{\sigma\mid\infty}n_\sigma\log\|s\|_\sigma\right).
\end{equation}
For a fractional ideal with the metrics below, the finite sum is
finite. Replacing $s$ by $xs$ adds zero: the finite norm sum for
$x\in K^\times$ is $-\log|N_{K/\Q}(x)|$, whereas the infinite
sum is $\log|N_{K/\Q}(x)|$. This proves independence of the rational
section with the indicated multiplicities.

For a fractional ideal $\mathfrak a\subset K$ with standard ambient
infinite metric $\|x\|_\sigma=|\sigma(x)|$, the section 1 gives
\begin{equation}\label{eq:ub-ideal-degree}
 \widehat{\deg}_K(\overline{\mathfrak a})
          =-[K:\Q]^{-1}\log N(\mathfrak a).
\end{equation}
In fact if $\mathfrak a_w=\pi_w^b\OO_{K_w}$, then
$\|1\|_w=q_w^b$. Thus a positive finite ideal exponent has negative
degree in this lattice convention. Degree of a vector bundle means
degree of its determinant, with the induced determinant metric.

\begin{theorem}[Arithmetic realization of the native hulls]
\label{thm:ub-realization}
Set
\begin{equation}\label{eq:ub-bundles}
 \mathfrak a_j=\prod_{w\mid p}\mathfrak p_w^{\eta_j},\qquad
 \mathcal E_j=\mathfrak a_j^{\oplus N_j}\subset K^{N_j},
\end{equation}
and give $\mathcal E_j$ the standard ambient Hermitian metric at
every infinite place. These are projective $\OO_K$-modules of rank
$N_j$. At each $w\mid p$ their local lattices identify with the
conjugate of $M_j^\rho$ over its first coefficient field. At all
other finite places they are the standard lattices. Moreover
\begin{equation}\label{eq:ub-degree-volume}
 \frac{\widehat{\deg}_K(\overline{\mathcal E_j})}{N_j}
       =-\frac{\eta_j}{e}\log p
       =V_j^B(M_j^\rho).
\end{equation}
\end{theorem}

\begin{proof}
A fractional ideal of the Dedekind domain $\OO_K$ is invertible,
hence projective of rank one. This applies also to negative $\eta_j$.
Since $K/\Q$ is Galois, all completions above $p$ have the same
ramification index $e$, residue degree $f$, and local degree $d$.
The product coordinates of $T_j$ identify its first-factor lattice
with $N_j$ copies of $\OO_{K_w}$. The depth $\eta_j$ of
\eqref{eq:ub-local-hull} in every component is precisely the depth
specified by \eqref{eq:ub-bundles}. Uniformizer units do not change
the lattice. Assigning the same exponent to every place over $p$
also makes $\mathfrak a_j$ Galois-stable.

If $g$ places lie above $p$, then $[K:\Q]=gef$ and
$N(\mathfrak p_w)=p^f$. Since
$\det\mathcal E_j=\mathfrak a_j^{N_j}$,
\eqref{eq:ub-ideal-degree} gives
\[
 \widehat{\deg}_K(\overline{\mathcal E_j})
 =-\frac{gfN_j\eta_j}{gef}\log p
 =-\frac{N_j\eta_j}{e}\log p.
\]
The rational section used here is the standard coordinate wedge,
of norm 1 at each infinite place. Thus the infinite terms are
exactly zero for the stated metrics, not omitted terms.

The local index calculation, valid for either sign of $\eta_j$,
gives $\mu_{B_j}(M_j^\rho)=p^{-fN_j\eta_j}$. Dividing its logarithm
by $D_j=efN_j$ proves \eqref{eq:ub-degree-volume}.
\end{proof}

Equivalently the contribution per rank is
\[
 \sum_{w\mid p}\frac{[K_w:\Q_p]}{[K:\Q]}
       \left(-\frac{\eta_j}{e}\log p\right)
       =-\frac{\eta_j}{e}\log p.
\]
Thus passing from one selected completion to all conjugate places
introduces no further degree or tensor-length factor.

\subsection{Weighted determinants, common twists, and descent}

\begin{proposition}\label{prop:ub-weighted-determinant}
Let $C$ be any positive integer divisible by every $hN_j$. The
metrized line
\begin{equation}\label{eq:ub-weighted-line}
 \overline{\mathcal L}
 =\bigotimes_{j=1}^h
       (\det\overline{\mathcal E_j})^{\otimes C/(hN_j)}
\end{equation}
is defined using positive integral tensor powers and satisfies
\begin{equation}\label{eq:ub-weighted-degree}
 \frac{\widehat{\deg}_K(\overline{\mathcal L})}{C}
       =\frac1h\sum_{j=1}^h V_j^B(M_j^\rho).
\end{equation}
Replacing every $\overline{\mathcal E_j}$ by
$\overline{\mathcal E_j}\otimes\overline Q$, for one metrized line
$\overline Q$, tensors \eqref{eq:ub-weighted-line} by
$\overline Q^{\otimes C}$.
\end{proposition}

\begin{proof}
The divisibility hypothesis makes all powers integral and positive.
Additivity of degree and \eqref{eq:ub-degree-volume} prove
\eqref{eq:ub-weighted-degree}. The canonical determinant isometry
\[
 \det(\overline{\mathcal E_j}\otimes\overline Q)
       =\det\overline{\mathcal E_j}\otimes\overline Q^{\otimes N_j}
\]
shows that the total exponent of $\overline Q$ is
$\sum_jN_jC/(hN_j)=C$.
\end{proof}

The reference determinant is present in this computation:
$\det\OO_K^{N_j}$ has the same standard metric and degree zero.
Each determinant in \eqref{eq:ub-weighted-line} can equivalently be
replaced by
$\det\mathcal E_j\otimes(\det\OO_K^{N_j})^{-1}$.
If covolumes are instead computed after restriction of scalars to
$\Q$, the field-discriminant and reference covolume factors appear
in both equal-rank terms and cancel in this relative determinant.

\begin{corollary}[Descent to the field of moduli]
\label{cor:ub-descent}
If $C$ is divisible by every $ehN_j$, put
\begin{equation}\label{eq:ub-descent-exponent}
 t=\frac{C}{eh}\sum_{j=1}^h\eta_j\in\Z.
\end{equation}
Then $\overline{\mathcal L}$ is the isometric pullback from $\Q$ of
the actual metrized line $p^t\Z\subset\Q$ with its ordinary ambient
absolute-value metric. In rational tensor powers,
$\overline{\mathfrak a_j}$ is the pullback of the arithmetic line
with finite exponent $\eta_j/e$ at $p$ and the standard infinite
metric.
\end{corollary}

\begin{proof}
Unique ideal factorization and the common ramification index give
\begin{equation}\label{eq:ub-descent-ideal}
 p\OO_K=\prod_{w\mid p}\mathfrak p_w^e,\qquad
 \mathfrak a_j^e=p^{\eta_j}\OO_K.
\end{equation}
Multiplication identifies the tensor metric on the left of the
second equality with the ambient metric on the right, since
absolute values multiply at every embedding. This proves the
rational-tensor-power statement.

In \eqref{eq:ub-weighted-line}, the exponent at each $w\mid p$ is
\[
 \sum_{j=1}^h\eta_jN_j\,\frac{C}{hN_j}
      =\frac Ch\sum_{j=1}^h\eta_j=et.
\]
There is no finite support elsewhere. Hence the actual lattice is
$p^t\OO_K$, and the multiplication identification preserves its
standard infinite metric. This is the stated pullback. In
particular its normalized degree is $-t\log p$.
\end{proof}

An explicit simultaneous choice is $C=ehd^h$, since $N_j=d^j$ and
$C/(hN_j)=e\,d^{h-j}\in\Z_{>0}$. Thus the descent concerns an actual
powered determinant object, not merely an equality of degrees.
Galois stability by itself would not have supplied unpowered
integral descent; \eqref{eq:ub-descent-ideal} is the additional step.

\subsection{The exact procession average}

\begin{proposition}\label{prop:ub-exact-average}
Let
\[
 R_\ell=\sum_{j=1}^h (2j^2\bmod\ell),
\]
where residues are taken in $\{1,\ldots,\ell-1\}$. Then
\begin{equation}\label{eq:ub-exact-average}
 \frac{\widehat{\deg}_K(\overline{\mathcal L})}{C\log p}
 =\frac{\ell+1}{12}-\frac{\ell+5}{60\ell}
        +\frac{2R_\ell}{\ell(\ell-1)}.
\end{equation}
Each summand $V_j^B(M_j^\rho)$ in
\eqref{eq:ub-weighted-degree} is strictly positive.
\end{proposition}

\begin{proof}
Since $\ell$ is odd prime and $1\le j\le h<\ell$, the indicated
residue is nonzero. The fractional part of $2j^2/\ell$ is therefore
at least $1/\ell>1/e$, which also verifies
$k_j=\lfloor2j^2/\ell+\kappa\rfloor$.
Elementary summation gives
\[
 \sum_{j=1}^h m_j=\frac{h(\ell+5)}4,\qquad
 \sum_{j=1}^h\left\lfloor\frac{2j^2}{\ell}\right\rfloor
       =\frac{\ell^2-1}{12}-\frac{R_\ell}{\ell}.
\]
For the second identity, sum $2j^2/\ell$ and subtract the
fractional parts. Substitution into the average
\[
 h^{-1}\sum_j(\kappa m_j-k_j)
\]
gives \eqref{eq:ub-exact-average}. Moreover $2j^2/\ell<j$, so $k_j\le j$
and
\[
 \kappa m_j-k_j\ge1-\frac{j+1}{e}>0.
\]
\end{proof}

As a numerical check at $\ell=43$, direct integer summation gives
$R_{43}=473$, $\sum_jk_j=164$, and $\sum_j\eta_j=-56508$.
The right side of \eqref{eq:ub-weighted-degree} evaluates to
$18836\log p/4515$. No distributional assertion about quadratic
residues is used here. In particular this positive degree is a
property of \eqref{eq:ub-bundles}, not a sign assertion about the
complete theta-pilot quantity in the source.

\subsection{The tensor-order correction}

\begin{proposition}\label{prop:ub-order-correction}
Let $\mu_{A_j}$ be additive Haar measure with $\mu_{A_j}(A_j)=1$.
For every measurable region $H\subset T_j$ of finite positive
volume,
\begin{equation}\label{eq:ub-order-correction}
 D_j^{-1}\log\mu_{A_j}(H)
   =D_j^{-1}\log\mu_{B_j}(H)
           +\frac{(m_j-1)\kappa}{2}\log p.
\end{equation}
\end{proposition}

\begin{proof}
Tameness gives the discriminant exponent $f(e-1)$ for
$\OO_E/\Z_p$. On a tensor integral basis of $A_j$, the Gram matrix
of the absolute algebra trace is the Kronecker product of the
factor trace Gram matrices. Its determinant therefore has
valuation
\[
 v_p\operatorname{disc}(A_j)=m_jd^{m_j-1}f(e-1).
\]
The product order $B_j$ has $N_j=d^{m_j-1}$ copies of $\OO_E$, so
\[
 v_p\operatorname{disc}(B_j)=d^{m_j-1}f(e-1).
\]
Under a finite-index change of lattice, the trace determinant
changes by the square of the index. Consequently
\[
 \log[B_j:A_j]
   =\frac{m_j-1}{2}d^{m_j-1}f(e-1)\log p.
\]
Finally $\mu_{A_j}=[B_j:A_j]\mu_{B_j}$. Dividing by
$D_j=efd^{m_j-1}$ proves the formula.
\end{proof}

The factor-weight prescription in
\cite[(9.10.3.1), (9.10.5.1)]{JoshiIII}, for this rational branch,
uses exponent $1/d$ on each factor. On tensor lattices this agrees
with $\mu_{A_j}^{1/D_j}$, since tensoring a lattice change in one
factor raises its determinant to the power $D_j/d$.
It therefore has the correction
\eqref{eq:ub-order-correction} relative to the maximal-order
reference. Agreeing on the factor weights alone does not remove
the $A_j/B_j$ index.

\subsection{All-slot scaling and transport of the infinite metrics}

\begin{proposition}\label{prop:ub-standard-scale}
Write $\mathcal E_j^\rho=\mathcal E_j$ for the realization
\eqref{eq:ub-bundles}. Replace all $m_j$ entries of the tuple
$(\tau_j,u,\ldots,u)$ specified in \eqref{eq:ub-native-points} by their
standard Bloch--Kato logarithms, retaining the same cohomology
classes, arrows, and ambient reference $B_j$. Then
\begin{equation}\label{eq:ub-standard-scale}
 \begin{aligned}
 M_j^{\rm std}&=p^{m_j}M_j^\rho
       =\beta^{\eta_j+em_j}B_j=\beta^{e k_j+m_j}B_j,\\
 V_j^B(M_j^{\rm std})&=V_j^B(M_j^\rho)-m_j\log p.
 \end{aligned}
\end{equation}
For the global realization, making this change only at the
specified rational prime $p$, while keeping the ambient infinite
metrics fixed, gives
\begin{equation}\label{eq:ub-standard-degree}
 \begin{gathered}
 \mathcal E_j^{\rm std}=p^{m_j}\mathcal E_j^\rho,\\
 \frac{\widehat{\deg}_K(\overline{\mathcal E_j^{\rm std}})
       -\widehat{\deg}_K(\overline{\mathcal E_j^\rho})}{N_j}
       =-m_j\log p.
 \end{gathered}
\end{equation}
The label-averaged change is $-(\ell+5)\log p/4$.
If instead the infinite metrics are transported along
multiplication by $p^{m_j}$, the map is an isometry of metrized
bundles and their degrees are equal.
\end{proposition}

\begin{proof}
The two coordinates differ by the rational scalar $p$ at every
entry, including all backgrounds. The permitted arrows are
$\Q_p$-linear, and the tensor map is multilinear. Thus every image
tensor is multiplied by $p^{m_j}$. Closure and $B_j$-span commute
with this invertible scalar. Its determinant on the
$D_j$-dimensional space has absolute value $p^{-m_jD_j}$,
which proves \eqref{eq:ub-standard-scale}.

At finite places away from $p$ the scalar is a unit, and at every
place above $p$ its uniformizer exponent is $em_j$. The degree
calculation in Theorem~\ref{thm:ub-realization} proves
\eqref{eq:ub-standard-degree}. The average of $m_j=j+1$ is
$(\ell+5)/4$.

For the transported metric, the new infinite norm on a vector $y$
is $p^{-m_j}\|y\|_\sigma$. The determinant norm is therefore
multiplied by $p^{-m_jN_j}$. Since
$\sum_{\sigma\mid\infty}n_\sigma=[K:\Q]$, its normalized infinite
degree contribution increases by $m_jN_j\log p$, cancelling the
finite change. This also follows from the resulting isometry.
\end{proof}

Thus scaling a finite coordinate with fixed references differs
from transporting a whole metrized global object. Neither
operation allows the active root alone to change scale while
the other source entries are silently left unchanged.

\subsection{An almost-everywhere reference test}

\begin{proposition}\label{prop:ub-ae-reference}
Fix one global member and one label $j$. Outside a finite set of
rational primes $p'$, take the declared background class
$\mathrm{Kum}(1+p')$ at the selected completion, its $m_j$-fold tuple,
and the native integral Kummer arrows. Relative to the maximal
order $B_{j,p'}$, its native hull is $B_{j,p'}$, whereas its
standard-coordinate hull is $(p')^{m_j}B_{j,p'}$.
With the reference unchanged, the sum of standard-coordinate
normalized local log-volumes is $-\infty$, and the volume product
is zero. This is not a finite-support fractional-ideal modification
of a fixed arithmetic bundle.
\end{proposition}

\begin{proof}
Exclude 2, the primes ramified in $K$, the bad primes of the curve,
and the other finitely many distinguished primes. The remaining
completion $E_{p'}/\Q_{p'}$ is unramified. The logarithm and
exponential on principal units give
\[
 \log\OO_{E_{p'}}^\times=p'\OO_{E_{p'}},\qquad
 I_{p'}=\OO_{E_{p'}},\qquad
 u_{p'}=(p')^{-1}\log(1+p')\in\Z_{p'}^\times.
\]
The residue-root-of-unity part has logarithm zero. The tensor
product of the unramified integral orders is finite etale over
$\Z_{p'}$ and is already its maximal product order.

Every integral automorphism of $I_{p'}$ preserves nonzero classes
modulo $p'I_{p'}$. Since the field is unramified, this is precisely
the unit condition. Each permitted image of $u_{p'}$ is therefore
a unit, and its repeated tensor has unit value in every field
component. The identity arrow is present. The $B_{j,p'}$-span of
the native images is consequently exactly $B_{j,p'}$.
The all-slot argument of Proposition~\ref{prop:ub-standard-scale}
gives the standard hull $(p')^{m_j}B_{j,p'}$, with normalized
log-volume $-m_j\log p'$, independently of the local degree.

Infinitely many rational primes remain, and
$\sum_{p'}m_j\log p'=+\infty$, since every summand is at least
$m_j\log2>0$. This proves the assertions about the sum and product.
A lattice in a fixed finite-dimensional $K$-space agrees with a
fixed integral reference at all but finitely many places. Nonzero
depth $m_j$ at infinitely many distinct rational primes violates
this condition. A finite contribution from fixed metrics at infinity
cannot cancel the divergence.
\end{proof}

Equivalently, the local scalars $(p')_{p'}$ do not form an idele:
their inverses are nonintegral at infinitely many primes. One
can replace the reference by $(p')^{m_j}B_{j,p'}$ at the same time,
obtaining local volume 1 again, but that specifies a different
restricted product and requires an explicit comparison with the
old reference.

\begin{remark}[The remaining identification]
The almost-everywhere test concerns the declared native branch,
not the complete Ind3 family. Likewise, the lattices assigned
away from $p$ in \eqref{eq:ub-bundles} and their standard infinite
metrics are part of our construction. They have not been identified
with the other components of the full theta-pilot output.
Theorems about these actual arithmetic bundles therefore supply
neither the particular membership in IUT III, Corollary~3.12,
nor the required comparison across Ind3. Establishing that
membership must retain the source classes, allowed arrows,
reference structures, weights, and the link to the same global
q-pilot. Equality of degrees of separately constructed objects
does not provide it.
\end{remark}
